\uuid{CWVs}
\titre{Structures non usuelles}
\auteur{hanani ; mzari}
\organisation{exo7}
\licence{CC BY-NC-SA 4.0}
\datecreate{2020-01-20}
\niveau{L1}
\module{Algèbre}
\chapitre{Espace vectoriel}
\sousChapitre{Définition, sous-espace}
\difficulte{3}
\theme{Espaces vectoriels, Structures non usuelles}
\exo7qcmid{lille-2.espaces-vectoriels.023}

\contenu{
\begin{question}
\qclassification{106.01}
\qtags{motcle=Structures non usuelles}
On considère $\Rr^2$ muni de l'addition et la multiplication par un réel suivantes :
$$(x,y)+(x',y')=(x+y',x'+y)\quad \mbox{et}\quad \lambda . (x,y) = (\lambda x, \lambda y).$$
Quelles sont les assertions vraies ?
\begin{answers}
\good{$E$ est stable par addition et par multiplication par un scalaire.}
\bad{L'addition est commutative.}
\bad{L'élément neutre pour l'addition est $(0,0)$.}
\bad{$E$ est un $\Rr$-espace vectoriel.}
\end{answers}
\begin{explanations} $E$ est stable par addition et multiplication par un scalaire. On voit que
$$(1,0)+(0,0)=(1,0)\quad \mbox{et}\quad (0,0)+(1,0)=(0,1)\neq (1,0).$$
On en déduit que l'addition n'est pas commutative et que $(0,0)$ n'est pas un élément neutre pour cette addition. En particulier, $E$ n'est pas un espace vectoriel.
\end{explanations}
\end{question}
}
